\chapter{Conclusion and Future Work}\label{ch:ch13}

This final chapter closes the thesis. It restates what APEX set out to do and
what was actually built, distils the engineering lessons that were paid for in
real incidents rather than borrowed from textbooks, and lays out a concrete
roadmap that turns each known limitation into a scoped piece of future work.
The intent throughout is to be specific and plain-spoken: the project names its
rough edges and explains exactly how it would fix each one.

\section{Summary of Contributions}

The central claim of this thesis is that mixed-format programming assessment
does not need a heavyweight platform to be taken seriously. APEX demonstrates
this by unifying three assessment formats that have historically lived in three
separate families of tools --- online judges, quiz engines, and document
hand-ins --- into a \emph{single} timed sitting with a \emph{single} grading
pipeline behind all three, packaged as a \emph{single} self-hostable artifact.
Concretely, the work delivered the following.

\begin{description}
  \item[A unified mixed-format examination portal.] Teachers author coding,
    multiple-choice (MCQ), and written (free-text) questions inside one exam,
    assign it to one or more classes, and set one timer. Students join a class
    by an eight-character invite code, sit the exam in a Monaco editor with
    live sample-test feedback, and review per-test results afterward. The
    system models fifteen first-class entities and exposes forty-seven feature
    pages across ten frontend feature areas, all behind exactly two roles ---
    \texttt{teacher} and \texttt{student} --- with no administrator role to
    complicate the trust model. The scale of the delivered surface is
    summarised in \cref{fig:loc} and \cref{fig:counts}.

  \item[One grading pipeline behind all three formats.] Rather than three
    graders bolted together, a single dispatch in the exam-submit handler
    routes each answer to \texttt{judge0.Grade} (coding),
    \texttt{judge0.GradeMCQ}, or \texttt{judge0.GradeWritten}, and every path
    ends by calling the same attempt aggregator, \texttt{maybeFinalizeAttempt}.
    Coding answers earn proportional partial credit
    (\(\texttt{passed}/\texttt{total}\times100\)); MCQ answers are graded by
    exact set-equality; written answers enter a manual-grading queue. The full
    path from the editor to a final score and a single notification is traced
    in \cref{fig:grading-seq}.

  \item[A sandboxed, fairness-corrected code execution engine.] Untrusted
    student code never runs on the application host. It is delegated over HTTP
    to Judge0~\cite{judge0,dosilovic2018judge0}, which executes it inside the
    \texttt{isolate} sandbox --- the same sandbox used at the International
    Olympiad in Informatics. APEX treats Judge0 as an \emph{executor}, not a
    \emph{judge}: it accepts both Judge0 status \texttt{3} (Accepted) and
    \texttt{4} (Wrong Answer) as ``the program ran'' and then re-compares the
    output itself through a nine-line \texttt{normalizeOutput} function applied
    symmetrically to the student's output and the expected output. Those nine
    lines eliminate the single most common false negative in I/O-diff grading:
    a correct answer failing on a trailing newline, a Windows CRLF line ending,
    or a stray trailing space.

  \item[A self-hostable one-artifact deployment.] The backend compiles to one
    statically linked Go binary (\texttt{CGO\_ENABLED=0}, built with
    \texttt{-trimpath -ldflags="-s -w"}) shipped in a
    \texttt{gcr.io/distroless/static-debian12:nonroot} image that runs as an
    unprivileged user with no shell. A production instance is therefore just
    that binary, a PostgreSQL~16 database, and an external Judge0 endpoint; the
    \texttt{docker-compose.yml} stack composes exactly four services ---
    \texttt{postgres}, a one-shot \texttt{migrate} job, \texttt{backend}, and
    an nginx-served \texttt{frontend} --- with Judge0 reached externally via
    \texttt{JUDGE0\_URL}. The deployment topology is shown in
    \cref{fig:deploy}.

  \item[Defence-in-depth authorization and a pre-ship security review.] Every
    request passes a four-layer lifecycle (CORS, authentication, role, handler)
    and a three-gate authorization model: HS256 JWT authentication, role
    checks, and per-handler resource-ownership checks. A security review run
    \emph{before} shipping surfaced four findings --- an IDOR on
    \texttt{/submissions/run}, a stored-XSS vector through Markdown
    \texttt{rehype-raw}, a weak fallback JWT secret, and HTML injection in
    outbound e-mail --- all fixed in a single commit prior to release. The
    three-gate model and threat model are detailed in \cref{ch:ch08} and
    depicted in \cref{fig:threat}.
\end{description}

\Cref{tab:objectives-review} maps the original project objectives onto these delivered
outcomes.

\begin{table}[htbp]
\centering
\caption{Project objectives and the outcome delivered against each.}
\label{tab:objectives-review}
\begin{tabularx}{\linewidth}{@{}L{4.3cm}X@{}}
\toprule
\textbf{Objective} & \textbf{Delivered outcome} \\
\midrule
Run coding, MCQ, and written questions in one timed exam &
Achieved --- one \texttt{ExamAttempt} per (student, exam), enforced by the
\texttt{idx\_user\_exam\_attempt} unique index; one submit endpoint creates all
answer rows atomically. \\
One grading pipeline behind all formats &
Achieved --- a single dispatch and a shared aggregator
(\texttt{maybeFinalizeAttempt}) with idempotent re-aggregation. \\
Safe execution of untrusted code &
Achieved --- delegated to Judge0/\texttt{isolate}; the host never executes
student code~\cite{judge0}. \\
Fair automatic grading &
Achieved --- symmetric output normalisation eliminates the trailing-whitespace
false negative. \\
Self-hostable with minimal moving parts &
Achieved --- one distroless binary, PostgreSQL, and an external Judge0; a
four-service compose stack. \\
Secure by construction &
Largely achieved --- three-gate authorization and a pre-ship review; several
hardening gaps remain and are named in \cref{sec:future}. \\
\bottomrule
\end{tabularx}
\end{table}

\section{Lessons Learned}

The most durable lessons of the project were not the ones that went smoothly.
Three in particular changed how the team writes code, and each is worth
recording because each was paid for with a real failure or a deliberate
trade-off rather than absorbed from a manual.

\subsection{The \texttt{is\_draft} Incident and Schema-Change Discipline}

The signature war story of the project --- a single
\texttt{gorm:"default:true"} struct tag that \texttt{AutoMigrate} applied as a
\emph{DDL} default to the already-populated \texttt{exams} table, silently
flipping every published exam out of every student's view with no error and no
stack trace --- is told in full in \cref{ch:ch05}. Its lesson is a standing
schema-change discipline: a boolean's zero value must be the safe,
non-destructive state; ownership of a populated-table change belongs in
explicit, idempotent SQL rather than a struct tag; and production runs with
\texttt{AutoMigrate} disabled, leaving a versioned \texttt{golang-migrate}
track~\cite{golangmigrate} as the sole authority over the schema.

That discipline --- audit every query when you add a filter column, define the
zero-value behaviour of every boolean before you add it, and verify row counts
after every migration --- was distilled into the \texttt{ORM\_RULES.md} rules
file, whose first rule is blunt: \textbf{never rely on \texttt{gorm:"default:X"}
for a DDL backfill on a populated table}. The general principle is old wisdom
about the gap between an ORM's convenient abstraction and the DDL it silently
emits~\cite{gorm,fowler2002patterns,kleppmann2017ddia}; APEX simply learned it
the expensive way. The two-track migration pipeline that resulted is shown in
\cref{fig:migration}.

\subsection{The Value of Idempotent Migrations}

The second lesson is a positive one that the first made unavoidable: every
migration in APEX is written to be safely re-runnable. The dev-mode legacy SQL
pass (\texttt{RunMigrations}) gates each block behind a \texttt{tableExists(...)}
check so that a first boot against an empty database does not spew ``relation
does not exist'' errors, and every column addition uses
\texttt{ADD COLUMN IF NOT EXISTS} before backfilling. The reminder scheduler
carries the same property into runtime: every reminder sweep dedups against a
dedicated timestamp column (\texttt{reminder\_sent\_at},
\texttt{email\_reminder1h\_sent\_at}, \texttt{email\_reminder\_start\_sent\_at}),
so a process restart mid-sweep can never re-send a reminder. Idempotency turned
out to be the cheapest possible form of operational safety: a migration or a
scheduled job that can be run twice without harm removes an entire category of
``did that already apply?'' anxiety from every deploy and every crash recovery.
It is the same principle the durable-queue future work below leans on ---
because grading recomputes results from the stored submission and test cases, it
too is idempotent, and idempotent work is safe to retry.

\subsection{Fire-and-Forget Versus Durable Work}

The third lesson is a trade-off the team made with its eyes open, and it defines
the largest single item of future work. Grading is dispatched as a bare
goroutine --- \texttt{go judge0.Grade(sub.ID)} --- so the HTTP handler can return
\texttt{201 Created} immediately; at the concurrency of a single classroom it
works, but a bare goroutine has no supervisor, no queue, and no retry, so a crash
mid-grade can strand a submission in \texttt{running} and leave the whole attempt
permanently unfinalised. The mechanism and its failure mode are traced in full in
\cref{ch:ch07} and documented openly in \texttt{README.md}. The lesson is not that
fire-and-forget was wrong for the scale it was built for; it is that
\emph{durability is a property you must design in deliberately}, and that the
right time to add a work queue is before you need to scale past one room, not
after a mid-exam crash. The design of that queue is the first item below, and
it is made tractable precisely by the idempotency lesson above.

\section{Future Work}\label{sec:future}

The remaining limitations of APEX are known, named, and prioritised. This
section presents each as a scoped roadmap item: the gap, its current
mitigation, and the concrete engineering that would close it.
\Cref{tab:roadmap} gives the priority ordering; the subsections that follow
give the detail.

\begin{table}[htbp]
\centering
\caption{Future-work roadmap: each remaining gap with its current mitigation and
planned fix.}
\label{tab:roadmap}
\begin{tabularx}{\linewidth}{@{}L{3.0cm}L{3.6cm}X@{}}
\toprule
\textbf{Gap} & \textbf{Current mitigation} & \textbf{Planned fix} \\
\midrule
No auth rate limiting &
bcrypt cost~10 only &
Per-IP and per-account limiter on auth routes; front with a WAF. \\
Fire-and-forget grading &
Works at classroom scale &
Durable job queue, bounded worker pool, restart recovery sweep. \\
JWT in \texttt{localStorage} &
No cookies, so no CSRF &
\texttt{httpOnly}~+~\texttt{SameSite} cookies; refresh-token model. \\
Minimal test coverage &
CI runs vet, lint, build &
Handler-level tests per protected route; grader unit tests. \\
Six execution languages &
Covers common cases &
Add languages via \texttt{LanguageMap}; push per-problem limits. \\
No plagiarism detection &
Hidden tests unshown &
MOSS / JPlag structural similarity over submissions. \\
Single-instance backend &
Sufficient for one class &
Stateless horizontal scaling plus metrics/tracing/logging. \\
\bottomrule
\end{tabularx}
\end{table}

\subsection{Hardening Authentication: Rate Limiting and a WAF}

The authentication endpoints have \emph{no} rate limiting. The only brake on
credential stuffing today is the deliberate cost of bcrypt at cost~10, and while
that meaningfully slows an offline or online guessing campaign, it does not stop
one. The cheapest and highest-value next mitigation is a per-IP \emph{and}
per-account rate limiter on \texttt{/auth/login}, \texttt{/auth/google}, and
\texttt{/auth/forgot-password}, implemented as a small Gin middleware backed by
an in-memory or Redis token bucket, coupled with exponential backoff on repeated
failures for a given account. For a public deployment, the README's standing
recommendation is to front the entire API with a web application firewall (WAF)
or an edge rate limiter so that abusive traffic is shed before it ever reaches
the application. The companion password-policy change --- raising the six-
character minimum in line with modern password-strength
guidance~\cite{nist80063b} --- is a one-line validation change that should ship
alongside the limiter. Both items sit at the top of the roadmap because,
together, they are the weakest part of the current authentication story, and
both align with the relevant guidance in the OWASP Top~Ten and the Application
Security Verification Standard~\cite{owasp-top10,owasp-asvs}.

\subsection{A Durable Grading Work Queue with Retries}

The largest architectural item is replacing the fire-and-forget grading of
\cref{ch:ch07} with a durable work queue. Concretely,
\texttt{go judge0.Grade(sub.ID)} would be replaced by enqueuing a job --- a row in a database-backed \texttt{grading\_jobs}
table, or a message on a Redis-backed broker such as \texttt{asynq} --- consumed
by a bounded pool of workers. Three properties follow from this design and
directly close the fire-and-forget gap:

\begin{itemize}
  \item \textbf{Crash recovery.} On startup, a recovery sweep re-picks any
    submission left in \texttt{running} (or any job left in-flight) and
    re-grades it. This is safe precisely because grading is
    \emph{idempotent}: \texttt{Grade} recomputes every \texttt{TestResult} from
    the stored \texttt{submission.Code} plus the problem's test cases and
    overwrites the previous results, so re-running it produces the same verdict.
    A stuck submission is no longer stuck forever; it is picked up on the next
    boot.
  \item \textbf{Bounded concurrency.} A worker pool caps the number of
    simultaneous Judge0 requests. Today a 200-student exam could spawn hundreds
    of concurrent POSTs and trip the public Judge0 CE rate limit; a pool with,
    say, a configurable maximum turns that burst into an orderly queue.
  \item \textbf{Retries with backpressure.} Transient Judge0 failures (the
    \texttt{StatusID = -1} sentinel today surfaces as a stuck or failed grade)
    become retriable jobs with a bounded retry count and a dead-letter path for
    permanent failures.
\end{itemize}

Because the attempt aggregator already re-sums the children on every call and
guards on ``no siblings pending or running,'' it needs no change to work with a
queue --- the last worker to finish still finds zero pending and finalises the
attempt exactly once, deduped by the \texttt{graded\_notified} flag. The queue
is an additive change, not a rewrite.

\subsection{\texttt{httpOnly} + \texttt{SameSite} Cookie Authentication}

The JWT is currently stored in the browser's \texttt{localStorage} under the key
\texttt{kernel-token} and sent as an \texttt{Authorization: Bearer} header.
Because no cookie is involved, the design has \emph{no} CSRF exposure by
construction --- but any successful cross-site scripting (XSS) would be able to
read \texttt{localStorage} and lift the token. The chosen trade was a simpler
deployment; the hardening path is to move the session token into an
\texttt{httpOnly}, \texttt{Secure}, \texttt{SameSite=Lax} (or \texttt{Strict})
cookie that JavaScript cannot read, which neutralises the token-theft-via-XSS
vector at the cost of reintroducing the need for CSRF defences (a
double-submit token or an origin check on state-changing requests). The same
change is the natural moment to introduce a proper refresh-token model with a
server-side revocation table, so that the current 24-hour non-revocable HS256
token can be invalidated on logout or compromise --- an improvement squarely in
line with established guidance on the safe use of JWTs and OAuth-style session
handling~\cite{rfc7519,owasp-jwt,rfc6749,nist80063b}.

\subsection{Real Test Coverage}

Test coverage is the most straightforward gap to close and, arguably, the most
valuable given how much of the system is security-sensitive. The backend
currently ships \emph{no} \texttt{*\_test.go} files, and the frontend ships a
single sanity test. Crucially, no pipeline work is required to add tests: the CI
already spins up a \texttt{postgres:16-alpine} service and runs
\texttt{go test ./... -race -count=1}, so the runners are in place and simply
have nothing to exercise. The priority targets are the ones where a regression
would be invisible and dangerous: a handler-level test per protected route
asserting that the three authorization gates reject the wrong role and the wrong
owner (an ownership-check regression is exactly the class of bug that the
pre-ship IDOR finding represents); unit tests over \texttt{normalizeOutput}
covering CRLF, lone CR, trailing whitespace, and trailing-newline cases; and
tests over \texttt{maybeFinalizeAttempt} asserting that skipped questions stay
in the denominator, that \texttt{pending\_review} answers are excluded from both
numerator and denominator, and that the ``Exam Graded'' notification fires
exactly once. Building this suite would move the project meaningfully up the
test pyramid described in \cref{fig:testpyramid}, following standard
software-engineering testing practice~\cite{sommerville-se,pressman-se}.

\subsection{More Execution Languages and Enforced Per-Problem Limits}

APEX supports six execution languages --- Python, JavaScript, TypeScript, Java,
C, and C++ --- resolved through the \texttt{LanguageMap} that maps a language
string to a Judge0 numeric identifier. Adding a language is a small, localised
change: a new entry in \texttt{LanguageMap} pointing at the corresponding Judge0
language id, and the rest of the pipeline is language-agnostic. Two adjacent
improvements belong here. First, the per-problem limits that already exist in
the schema --- \texttt{Problem.TimeLimitMs} (default \texttt{2000}) and
\texttt{Problem.MemoryLimitKb} (default \texttt{262144}, i.e.\ 256~MB) --- are
\emph{not} currently pushed down to Judge0; execution is bounded by Judge0's own
defaults instead. Wiring \texttt{cpu\_time\_limit} and \texttt{memory\_limit}
into the \texttt{RunCode} request body is a small change that would let a teacher
who sets a one-second limit actually have it enforced. Second, the grader
currently sends and receives plain text (\texttt{base64\_encoded=false}); moving
to base64 transport would make the pipeline robust to exotic byte sequences in
inputs and outputs.

\subsection{Plagiarism Detection}

APEX has no plagiarism detection today. Because the platform already stores every
submission's source code alongside its language and the exam it belongs to, it
is well positioned to run structural-similarity analysis across a cohort's
submissions after an exam closes. The obvious approach is to integrate a
document-fingerprinting detector in the style of MOSS~\cite{moss} or a
token-based structural comparator such as JPlag~\cite{prechelt2002jplag}, both of
which compare programs by structure rather than surface text and are therefore
robust to renamed variables and reordered statements. The feature would surface a
ranked list of suspiciously similar submission pairs to the teacher as a review
aid --- explicitly a signal for human judgement, not an automatic accusation.
This pairs naturally with lighter-weight proctoring signals (tab-blur,
fullscreen-exit, and paste events captured client-side) as a future integrity
layer.

\subsection{Horizontal Scaling and Observability}

The backend runs today as a single instance, which is entirely sufficient for
one classroom but is a ceiling for wider deployment. The good news is that the
application is already close to horizontally scalable: it is a stateless Go
process (all state lives in PostgreSQL) that follows the twelve-factor
principles of configuration-through-environment and disposable
processes~\cite{twelvefactor}. The two things standing between the current design
and running \(N\) replicas behind a load balancer are both addressed by the work
above --- the in-process reminder scheduler would need a leader-election lock so
that only one replica runs the sweeps (or it would move into the same durable
job system as grading), and grading would move from in-process goroutines to the
shared work queue so that any worker can grade any submission. Alongside scaling,
the project needs real observability: structured request logging, Prometheus-
style metrics on grading latency and queue depth, and distributed tracing across
the SPA~\(\rightarrow\)~backend~\(\rightarrow\)~Judge0 path. Today the system has
liveness (\texttt{GET /healthz}) and a self-probing \texttt{--healthcheck} flag
but no metrics or tracing; adding them is what turns ``it seems to work'' into
``we can see it working, and we get paged before the student does.''

\section{Closing Statement}

APEX set out to build the smallest system that still takes mixed-format
classroom assessment seriously, and it succeeded on its own terms: three
assessment formats in one timed sitting, one grading pipeline behind them, and
one self-hostable artifact of a Go binary, PostgreSQL, and an external Judge0
sandbox. Along the way the project learned --- expensively, from a single struct
tag that hid every exam from every student --- that the boundary between an ORM's
convenience and the DDL it silently emits is exactly where data integrity is won
or lost, and it built a schema-change discipline and an idempotent,
crash-resilient runtime around that lesson. The limitations that remain are not
hidden: no auth rate limiting, fire-and-forget grading, a token in
\texttt{localStorage}, thin test coverage, and no plagiarism detection. Each is
named in this thesis, each has a concrete fix on the roadmap above, and each was
a deliberate trade for a system that a single instructor can actually stand up
and run. That honesty is the point. The most useful thing a graduation project
can leave behind is not the illusion of a finished product but a working system
whose author knows exactly where its edges are and exactly how to extend them ---
and that is what APEX is.
